\chapter{Amputation}

\section*{Introduction \& Clinical Indications}
Amputation is the surgical removal of a limb, digit, or part of a limb when the tissue is nonviable, severely infected, or irreparably damaged and cannot be salvaged through reconstructive means. This procedure is most commonly performed for advanced peripheral arterial disease (PAD) leading to critical limb ischemia and gangrene, severe necrotizing infections such as osteomyelitis or fasciitis, or devastating trauma. While often viewed as a last resort, amputation can be a life-saving procedure, preventing the spread of infection or systemic toxicity, and can significantly improve a patient's quality of life by eliminating chronic pain and allowing for prosthetic rehabilitation. The primary goal is to create a functional, pain-free residual limb (stump) that can support a prosthesis and facilitate mobility, thereby restoring a degree of independence and improving overall well-being.

\begin{itemize}
  \item Limb removal
  \item Lower-extremity amputation (LEA)
  \item Upper-extremity amputation (UEA)
  \item Transfemoral amputation (Above-knee amputation, AKA)
  \item Transtibial amputation (Below-knee amputation, BKA)
  \item Disarticulation (e.g., knee disarticulation, hip disarticulation)
\end{itemize}

The fundamental principle of amputation is the complete excision of nonviable, infected, or irreparably damaged tissue at a level that possesses adequate blood supply to ensure primary wound healing. The primary goals are multifaceted: to remove the source of disease (e.g., gangrene, necrotizing infection, unresectable tumor), alleviate intractable pain, prevent systemic complications such as sepsis, and create a functional residual limb (stump) that is well-padded, non-tender, and amenable to prosthetic fitting. This involves careful preoperative assessment of the vascularity of the proposed amputation level, often utilizing clinical examination, Doppler studies, transcutaneous oxygen measurements, or angiography, to predict healing potential.

Surgical technique focuses on preserving maximal limb length while ensuring healthy tissue margins and optimal soft tissue coverage. The bone is divided cleanly and smoothed to prevent sharp edges that could cause pressure points. Muscles are meticulously shaped and often reattached to the bone (myodesis) or to opposing muscle groups (myoplasty) to provide padding, maintain muscle balance, and improve prosthetic control. Major arteries and veins are ligated to achieve hemostasis, and nerves are transected under tension (traction neurectomy) to allow them to retract into soft tissue, minimizing the risk of painful neuromas. The skin flaps are designed to provide durable coverage without tension, with the scar ideally placed away from weight-bearing areas to optimize prosthetic use and patient comfort.

Amputation is one of the oldest surgical procedures, with archaeological evidence dating back to prehistoric times. Ancient Egyptian and Greek texts, including those by Hippocrates (c. 460--370 BC), describe limb removal for gangrene, though survival rates were extremely low due to uncontrolled hemorrhage and infection. A significant advancement occurred in the 16th century with Ambroise Par\'e (1510--1590), a French barber-surgeon, who reintroduced the use of ligatures to control bleeding from major vessels, dramatically improving patient survival compared to the previously common practice of cauterization.

The 19th century brought further revolutions with the advent of anesthesia (ether, chloroform) and antisepsis (Joseph Lister's carbolic acid), transforming amputation from a brutal, rapid procedure into a more controlled operation with reduced pain and infection rates. The American Civil War and the World Wars spurred significant advancements in surgical techniques, prosthetic design, and rehabilitation, as large numbers of soldiers required amputations. Surgeons like Stephen Smith in the U.S. and Robert Jones in the UK refined techniques for creating functional stumps. In the latter half of the 20th century, the rise of vascular surgery and microsurgery shifted the focus towards limb salvage, making amputation a procedure of last resort, often preceded by attempts at revascularization. Modern amputation techniques prioritize functional outcomes, pain management, and early prosthetic fitting, integrating a multidisciplinary approach to patient care.

Amputation is indicated for a variety of severe conditions where limb salvage is impossible or has failed, or when the limb poses a threat to the patient's life or overall health. Specific indications include:

\begin{itemize}
  \item Critical limb ischemia (CLI) with nonviable tissue, such as wet or dry gangrene, severe rest pain, or non-healing ulcers, especially in patients with diabetes mellitus or peripheral arterial disease (e.g., Rutherford Classification categories 4-6).
  \item Severe, uncontrolled infection of the limb, including necrotizing fasciitis, osteomyelitis, or extensive cellulitis that is unresponsive to antibiotics and debridement, particularly when systemic sepsis is present.
  \item Irreparable traumatic injury, such as severe crush injuries, degloving injuries, avulsions, or blast injuries, where the limb is mangled beyond reconstruction or functional recovery is deemed impossible (e.g., high MESS score).
  \item Malignant tumors of bone or soft tissue (e.g., osteosarcoma, chondrosarcoma, rhabdomyosarcoma) that are extensive, aggressive, or unresponsive to chemotherapy and radiation, and cannot be adequately resected with limb-sparing surgery.
  \item Severe frostbite or thermal burns causing full-thickness tissue necrosis and gangrene that extends beyond the superficial layers and involves deeper structures.
  \item Congenital deformities or anomalies that result in a non-functional limb, severe pain, or recurrent infections, and cannot be corrected by other surgical means.
  \item Intractable pain or chronic infection in a non-functional limb, even without immediate life threat, when it significantly impairs the patient's quality of life and other palliative measures have failed.
\end{itemize}

Amputation can serve as both a primary and a secondary definitive treatment, depending on the clinical context and the urgency of the situation. It is considered a primary definitive treatment in emergent situations where immediate limb removal is necessary to save the patient's life or prevent catastrophic complications. Examples include severe, unreconstructable traumatic injuries with massive tissue destruction and hemorrhage, or fulminant necrotizing infections causing rapid systemic sepsis and multi-organ failure. In these critical cases, the decision for amputation is often made rapidly to stabilize the patient and remove the source of life-threatening pathology.

More commonly, amputation functions as a secondary or salvage procedure, performed after attempts at limb salvage have failed or are deemed futile. This is particularly true for patients with critical limb ischemia due to peripheral arterial disease or diabetic foot complications, where revascularization procedures (e.g., bypass surgery, angioplasty) have been unsuccessful, or the extent of tissue loss and infection is too advanced for salvage. In such scenarios, amputation becomes the most appropriate option to resolve chronic pain, eliminate infection, and allow for rehabilitation, ultimately improving the patient's functional status and quality of life. The decision to proceed with amputation is typically made after a thorough multidisciplinary evaluation, weighing the risks and benefits of continued limb salvage efforts against the potential for improved outcomes with amputation.

\section*{Relevant Surgical Anatomy}
A comprehensive understanding of surgical anatomy is paramount for successful amputation, focusing on the vascularity, innervation, musculature, and skeletal structure at the proposed level of transection. For lower extremity amputations, which constitute the majority, the major arterial supply originates from the aorta, descending through the common iliac, external iliac, common femoral, superficial femoral, and popliteal arteries, before branching into the anterior tibial, posterior tibial, and peroneal arteries. Adequate blood supply at the amputation level is critical for wound healing, often assessed by clinical signs, Doppler studies, or transcutaneous oxygen measurements. Venous drainage largely parallels the arterial system, with deep veins accompanying the arteries and superficial veins (e.g., great and small saphenous) draining into the deep system.

Nerve management is crucial to prevent painful neuromas. Major nerves such as the sciatic nerve (which branches into the tibial and common peroneal nerves) and the femoral nerve must be identified. These are typically transected under tension (traction neurectomy) to allow their ends to retract into soft tissue, minimizing irritation. Lymphatic drainage generally follows the venous pathways, with regional lymph nodes (e.g., inguinal nodes for the lower limb) being relevant in cases of infection or malignancy. Skeletal landmarks guide the level of bone transection, aiming to preserve maximal functional length while ensuring healthy margins. Myocutaneous flaps are meticulously designed to provide durable, well-padded coverage over the bone end, with muscle groups often reattached (myodesis or myoplasty) to stabilize the residual limb and facilitate prosthetic control.

\section*{Preoperative Workup \& Preparation}
Comprehensive preoperative preparation is essential to optimize patient outcomes and minimize complications, often involving a multidisciplinary approach. This includes a thorough medical evaluation to assess and optimize comorbidities, detailed imaging, and patient counseling.

\begin{itemize}
  \item \textbf{Medical Optimization:}
    \begin{itemize}
      \item \textbf{Glycemic Control:} For diabetic patients, perioperative glucose management is individualized; a universal HbA1c target should not delay urgent source control or lifesaving amputation.
      \item \textbf{Nutritional Assessment:} Malnourished patients are at higher risk for poor wound healing and infection. Nutritional support, sometimes including parenteral or enteral feeding, may be initiated if time permits.
      \item \textbf{Cardiovascular and Pulmonary Assessment:} Patients often have significant comorbidities. Cardiac clearance (ECG, stress tests, echocardiogram) and pulmonary function tests may be required to assess fitness for anesthesia and surgery, especially given the high prevalence of PAD.
      \item \textbf{Renal Function:} Assessment of kidney function (creatinine, GFR) is important for medication dosing and fluid management.
      \item \textbf{Coagulation Profile:} Blood tests to assess clotting ability (PT/INR, PTT, platelet count) are performed, and anticoagulant medications may need to be adjusted or temporarily discontinued.
    \end{itemize}
  \item \textbf{Infection Management:} If active infection is present, broad-spectrum intravenous antibiotics are typically started preoperatively and continued postoperatively, guided by culture results from wound swabs or tissue biopsies.
  \item \textbf{Vascular Assessment:} Non-invasive studies (e.g., ankle-brachial index, toe-brachial index, Doppler ultrasound, transcutaneous oxygen measurements) and sometimes angiography are performed to precisely determine the level of adequate blood supply for healing, guiding the surgical level selection.
  \item \textbf{Patient Counseling and Consent:} Detailed discussion with the patient and family regarding the procedure, potential risks, expected outcomes, and the extensive rehabilitation process is crucial. Psychological support and pre-amputation counseling are often beneficial to prepare the patient for the physical and emotional challenges.
  \item \textbf{NPO Status:} Patients are instructed to be NPO (nil per os - nothing by mouth) for a specified period (typically 6-8 hours) before surgery to prevent aspiration during anesthesia.
  \item \textbf{Prophylactic Antibiotics:} Intravenous antibiotics, usually a first or second-generation cephalosporin, are administered shortly before the incision to reduce the risk of surgical site infection.
\end{itemize}

While there are few absolute contraindications to amputation when a limb is nonviable and poses an immediate threat to life, certain conditions warrant careful consideration and optimization.

\textbf{Situations requiring clarification or optimization:}
\begin{itemize}
  \item \textbf{Severe instability or coagulopathy:} These require resuscitation and correction when feasible, but urgent amputation may still be necessary for source control or a nonviable limb.
  \item \textbf{Severe Uncorrectable Coagulopathy:} An extreme bleeding disorder that cannot be corrected with blood products or medications (e.g., Factor VIIa, desmopressin) poses an unacceptable risk of hemorrhage during and after surgery.
  \item \textbf{Patient Refusal:} A mentally competent patient's informed refusal of the procedure is an absolute contraindication, though extensive counseling and psychological support should be provided to ensure the decision is fully informed.
\end{itemize}

\textbf{Relative Contraindications and Precautions:}
\begin{itemize}
  \item \textbf{Severe Cardiopulmonary Disease:} Patients with severe heart failure (NYHA Class IV), unstable angina, or advanced respiratory disease (e.g., severe COPD) require intensive preoperative optimization and careful anesthetic management. The risks of surgery must be meticulously weighed against the benefits, and often a lower-level amputation is preferred if feasible.
  \item \textbf{Profound Malnutrition:} Poor nutritional status (e.g., albumin <2.5 g/dL) significantly impairs wound healing and increases infection risk. Nutritional support should be initiated preoperatively if time permits, and surgery may be delayed.
  \item \textbf{Inadequate Blood Supply at Planned Level:} If preoperative assessment indicates insufficient vascularity at the chosen amputation level, a higher level of amputation or further revascularization attempts may be necessary to ensure healing. Intraoperative assessment of bleeding from skin edges is crucial.
  \item \textbf{Active Infection Proximal to Planned Level:} While the goal is to remove infected tissue, if there is active, spreading infection proximal to the planned amputation site, a higher level may be required, or the infection needs to be controlled first with antibiotics and debridement.
  \item \textbf{Psychological Instability or Lack of Support System:} Patients with severe depression, cognitive impairment, or without adequate social support may struggle significantly with the rehabilitation process. Preoperative psychological assessment and support planning are crucial.
  \item \textbf{Marking the Level:} Always mark the planned amputation level clearly on the skin preoperatively, ideally with the patient awake and involved in the decision, and re-evaluate viability intraoperatively based on tissue perfusion.
  \item \textbf{Avoid Overly Aggressive Amputation:} While ensuring a healthy healing level is paramount, unnecessary removal of viable limb length should be avoided to maximize functional potential and prosthetic options.
  \item \textbf{Contralateral Limb Assessment:} Thoroughly assess the contralateral limb for signs of vascular disease, neuropathy, or pre-ulcerative lesions, as these patients are at high risk for future problems in the remaining limb.
\end{itemize}

\section*{Surgical Technique \& Procedure}
Amputation is performed in an operating room under general or regional anesthesia (e.g., spinal or epidural block), often combined with peripheral nerve blocks for enhanced postoperative pain control. Patient positioning depends on the level of amputation; for lower extremity amputations, the patient is typically supine with the affected limb prepped and draped to allow full access. A tourniquet may be used for upper extremity or distal lower extremity amputations to minimize blood loss, though it is generally avoided in cases of critical limb ischemia.

The surgical procedure involves several critical steps, meticulously executed to create a functional, well-healed residual limb:

\begin{itemize}
  \item \textbf{Incision and Flap Design:} The surgeon carefully plans and incises the skin and subcutaneous tissue, creating anterior and posterior flaps (or other designs like sagittal or circular) that will provide adequate coverage for the residual limb without tension. The length and shape of these flaps are crucial for optimal healing and prosthetic fit, often designed to place the scar away from weight-bearing areas.
  \item \textbf{Muscle Division and Shaping:} Muscles are divided at the appropriate level. For transtibial amputations, the gastrocnemius and soleus muscles are often shaped to create a posterior flap that provides padding over the tibia. In transfemoral amputations, muscles are often beveled and then secured to the bone (myodesis) or to opposing muscle groups (myoplasty) to stabilize the stump, maintain muscle balance, and improve prosthetic control.
  \item \textbf{Vessel Ligation:} Major arteries and veins supplying the distal limb are identified, carefully ligated with non-absorbable sutures, and then divided to ensure complete hemostasis. Smaller vessels are meticulously cauterized.
  \item \textbf{Nerve Management:} Peripheral nerves are identified, gently pulled distally, sharply transected with a scalpel, and allowed to retract proximally into the soft tissues. This "traction neurectomy" technique aims to minimize the formation of painful neuromas at the stump surface.
  \item \textbf{Bone Transection and Smoothing:} The bone (e.g., tibia and fibula for BKA, femur for AKA) is precisely transected using an oscillating saw at the predetermined level. The sharp edges and corners of the bone are then carefully beveled and rasped smooth to prevent pressure points, skin breakdown, and potential osteomyelitis within the prosthesis.
  \item \textbf{Hemostasis and Wound Closure:} Thorough hemostasis is achieved across all tissue planes. A drain (e.g., Jackson-Pratt or Blake drain) may be placed to prevent hematoma formation and reduce seroma risk. The muscle flaps are approximated, and the skin flaps are closed primarily with sutures or staples, ensuring no tension. In cases of significant infection or contamination, the wound may be left open for delayed primary closure or secondary intention healing.
\end{itemize}

A sterile dressing is applied, and often a soft or rigid dressing (e.g., plaster cast or prefabricated protector) is used to protect the stump, control edema, and begin shaping the residual limb.

\section*{Complications \& Risks}
Amputation, like any major surgical procedure, carries inherent risks, which can be broadly categorized into general surgical risks and procedure-specific complications.

\textbf{General Surgical Risks:}
\begin{itemize}
  \item \textbf{Bleeding and Hematoma:} Excessive blood loss during surgery or accumulation of blood (hematoma) in the stump postoperatively, which may require reoperation, blood transfusion, or lead to wound dehiscence.
  \item \textbf{Infection:} Surgical site infection (SSI) of the stump, which can range from superficial cellulitis to deep abscess or osteomyelitis, potentially leading to wound dehiscence, delayed healing, or the need for re-amputation.
  \item \textbf{Anesthesia Complications:} Risks associated with general or regional anesthesia, including adverse drug reactions, respiratory depression, cardiac events (e.g., myocardial infarction, arrhythmia), stroke, or nerve damage.
  \item \textbf{Deep Vein Thrombosis (DVT) and Pulmonary Embolism (PE):} Formation of blood clots in the deep veins, which can dislodge and travel to the lungs, causing a life-threatening PE. Prophylaxis is essential.
  \item \textbf{Cardiovascular Complications:} Given the high prevalence of cardiovascular disease in patients undergoing amputation, risks of heart attack, stroke, or exacerbation of heart failure are significant.
  \item \textbf{Renal Dysfunction:} Acute kidney injury can occur, especially in patients with pre-existing renal impairment or those experiencing significant blood loss, sepsis, or contrast-induced nephropathy from preoperative imaging.
  \item \textbf{Pneumonia:} Postoperative respiratory complications, particularly in elderly or debilitated patients, or those with pre-existing lung disease.
\end{itemize}

\textbf{Procedure-Specific Risks:}
\begin{itemize}
  \item \textbf{Phantom Limb Pain (PLP):} A common and often debilitating sensation of pain originating from the amputated limb, which can be difficult to manage and requires multimodal treatment strategies.
  \item \textbf{Stump Pain:} Pain localized to the residual limb, which can be due to neuromas, bone spurs, poor prosthetic fit, ongoing ischemia, or infection.
  \item \textbf{Stump Neuroma:} A painful, benign growth of nerve tissue at the transected nerve ending, which can cause localized tenderness and radiating pain, sometimes requiring surgical excision.
  \item \textbf{Wound Dehiscence or Non-Healing:} Failure of the surgical wound to close or heal properly, often due to infection, poor blood supply, excessive tension, or patient comorbidities, potentially requiring further surgery or revision.
  \item \textbf{Re-amputation at a Higher Level:} If the initial amputation site fails to heal due to inadequate blood supply or uncontrolled infection, a second surgery at a more proximal level may be necessary, leading to greater functional loss.
  \item \textbf{Contractures:} Shortening of muscles and tendons around the joint proximal to the amputation (e.g., hip flexion contracture after AKA, knee flexion contracture after BKA), limiting range of motion and hindering prosthetic fitting and ambulation.
  \item \textbf{Heterotopic Ossification:} Abnormal bone formation in the soft tissues of the residual limb, which can cause pain, limit range of motion, and interfere with prosthetic use.
  \item \textbf{Psychological Distress and Depression:} Significant emotional and psychological challenges, including grief, body image issues, anxiety, and depression, are common after limb loss and require dedicated support.
  \item \textbf{Contralateral Limb Complications:} Patients undergoing amputation, especially for vascular disease, are at high risk for developing similar problems in their remaining limb due to systemic disease progression.
  \item \textbf{Mortality:} Amputation, particularly for critical limb ischemia in elderly patients with multiple comorbidities, carries a significant risk of perioperative and long-term mortality, primarily due to underlying systemic disease.
\end{itemize}

\section*{Postoperative Care \& Recovery}
Immediately after the amputation, the patient is transferred to the Post-Anesthesia Care Unit (PACU) for close monitoring of vital signs, pain levels, and the surgical site. Pain management is a critical priority, often involving a combination of intravenous analgesics, regional nerve blocks, and patient-controlled analgesia (PCA) to ensure comfort and facilitate early mobilization. The stump dressing is regularly checked for bleeding or signs of infection, and drains are monitored for output. Early mobilization, often starting within 24-48 hours with assistance, is encouraged to prevent complications like deep vein thrombosis, pneumonia, and muscle deconditioning.

Over the first 1-2 weeks, the focus shifts to meticulous wound care, continued pain management, and the initiation of intensive rehabilitation. The surgical dressing is changed, and the wound is inspected for signs of healing or complications such as infection or dehiscence. Edema control is crucial, often managed with elastic bandages or stump shrinkers to reduce swelling and begin shaping the residual limb for future prosthetic fitting. Physical therapy commences with range-of-motion exercises for the proximal joints to prevent contractures, strengthening exercises for the core and remaining limb, and instruction on proper positioning. Occupational therapy addresses activities of daily living (ADLs) and adaptive strategies. Psychological support is often integrated early to help patients cope with limb loss and body image changes. Discharge planning includes arrangements for home care, outpatient physical therapy, and prosthetic consultation, typically occurring within 1-2 weeks for uncomplicated cases. Long-term recovery involves months to years of prosthetic training and adaptation.

During the amputation procedure, the surgeon meticulously documents the gross surgical findings, including the extent of nonviable or infected tissue, the quality of the remaining healthy tissue at the chosen amputation level, the appearance of major vessels and nerves, and any unexpected anatomical variations. This intraoperative assessment guides the final level of bone and soft tissue transection. Following removal, the amputated specimen is sent to pathology for detailed macroscopic and microscopic examination. The pathology report is crucial as it confirms the underlying diagnosis, such as the extent of atherosclerotic disease, the presence and type of necrotizing infection (e.g., osteomyelitis, fasciitis, specific bacterial species), or the specific malignancy (e.g., osteosarcoma, chondrosarcoma) with assessment of tumor margins. It also provides vital information on the health of the tissue at the surgical margins, confirming that all diseased or nonviable tissue was removed and that the remaining tissue is healthy, which is predictive of wound healing.

A normal, successful postoperative course following an amputation is characterized by several key indicators that signify proper healing and readiness for rehabilitation. The patient should experience effective control of postoperative pain, allowing for participation in early rehabilitation activities. The surgical incision should heal cleanly without signs of infection (e.g., increasing redness, warmth, swelling, purulent drainage) or dehiscence (wound separation). The skin flaps must remain viable, with no evidence of necrosis. Crucially, the underlying condition that necessitated the amputation (e.g., critical limb ischemia, active infection, severe pain from gangrene) should be resolved, with no ongoing signs of sepsis or spreading infection. The residual limb should gradually form a well-contoured, non-tender, and durable stump, free from excessive edema, painful neuromas, or bony prominences, which is essential for comfortable and effective prosthetic fitting. Furthermore, the joint proximal to the amputation (e.g., knee for BKA, hip for AKA) should maintain a good range of motion, free from contractures, which is crucial for prosthetic use and functional mobility.

Post-amputation care involves a structured, multidisciplinary follow-up plan to monitor healing, manage complications, and facilitate comprehensive rehabilitation.

\begin{itemize}
  \item \textbf{Regular Postoperative Visits:} Scheduled appointments with the surgical team (typically at 1-2 weeks, then monthly) to monitor wound healing, assess for signs of infection or dehiscence, and manage pain.
  \item \textbf{Suture/Staple Removal:} Typically occurs 2-4 weeks post-surgery, once the wound edges have adequately coapted and healing is progressing well.
  \item \textbf{Stump Shrinker/Compression Therapy:} Application of elastic bandages or custom-fitted stump shrinkers is initiated early to reduce edema, shape the residual limb into a conical or cylindrical form, and prepare it for prosthetic fitting.
  \item \textbf{Physical and Occupational Therapy:} Intensive, ongoing rehabilitation is crucial. This includes exercises to strengthen core muscles, improve balance, prevent contractures, and eventually gait training with a prosthesis. Occupational therapy focuses on adapting activities of daily living (ADLs) and promoting independence.
  \item \textbf{Prosthetic Evaluation and Fitting:} Once the stump is well-healed, stable, and has achieved a mature shape (typically 4-12 weeks post-op), the patient will be evaluated by a prosthetist for custom prosthetic limb fabrication and fitting. This process involves multiple adjustments and training sessions.
  \item \textbf{Pain Management Clinic Referral:} For persistent or severe phantom limb pain or stump pain, referral to a specialized pain management clinic may be necessary for advanced therapies, including medications, nerve blocks, or neuromodulation.
  \item \textbf{Psychological Support:} Ongoing counseling, support groups, or psychiatric evaluation may be recommended to address the emotional and psychological impact of limb loss, including grief, depression, and body image issues.
  \item \textbf{Contralateral Limb Surveillance:} For patients with vascular disease, regular monitoring of the remaining limb for signs of ischemia, neuropathy, or new wounds is vital to prevent future amputations.
\end{itemize}
Patients are educated on warning signs such as fever, increasing redness, swelling, severe pain, or foul-smelling drainage from the stump, and instructed to contact their healthcare provider immediately if these occur.

\section*{Outcomes \& Prognosis}
Amputation, particularly of the lower extremity, is a significant public health concern globally, with incidence rates varying by region and population demographics. In developed countries, the incidence of major lower extremity amputation (LEA) ranges from 12 to 50 per 100,000 population per year. The vast majority (80-90\%) of non-traumatic amputations are performed due to complications of peripheral arterial disease (PAD) and diabetes mellitus, often leading to critical limb ischemia, non-healing ulcers, and gangrene. The incidence of LEA is disproportionately higher in older adults, with rates increasing significantly after age 60, and is generally higher in males than females. Racial and socioeconomic disparities also exist, with higher rates observed in certain minority populations and underserved communities, often reflecting unequal access to preventative care and revascularization. Traumatic amputations account for a smaller percentage (5-10\%) but are more common in younger, working-age individuals. Amputations for malignancy or congenital deformities are relatively rare. Despite advancements in limb salvage techniques, the mortality rate following major LEA remains substantial, with 1-year mortality rates ranging from 20-30\% and 5-year mortality rates often exceeding 50\%, primarily due to the severe underlying comorbidities (cardiovascular disease, renal failure, diabetes) prevalent in this patient population.

The outcomes and prognosis following amputation are highly variable, influenced by the underlying indication, the level of amputation, patient comorbidities, and the quality of postoperative rehabilitation. For patients undergoing amputation for critical limb ischemia, wound healing rates for transtibial amputations are generally 70-85\%, while transfemoral amputations have higher healing rates (85-95\%) due to better proximal blood supply, but come with greater functional impairment. Functional outcomes are generally better for lower-level amputations (e.g., transtibial vs. transfemoral), as they allow for more efficient prosthetic use and ambulation. However, even with successful prosthetic fitting, many patients experience reduced mobility and require assistive devices, with only about 50-70\% of transtibial amputees and 20-40\% of transfemoral amputees achieving independent ambulation with a prosthesis.

Survival outcomes are largely dictated by the patient's comorbidities. As noted, 5-year mortality rates are high, often exceeding 50\% for patients with diabetes and PAD, reflecting the systemic nature of their disease. Recurrence is a significant concern; patients who undergo amputation for vascular disease have a high risk of developing similar problems in the contralateral limb, with a 5-year contralateral amputation rate of 15-25\%. Prognostic factors for better outcomes include younger age, fewer comorbidities, a lower level of amputation, good nutritional status, strong social support, and adherence to a rigorous rehabilitation program. While amputation can be life-saving and significantly improve quality of life by eliminating pain and infection, it represents a profound life change requiring extensive physical and psychological adaptation.

\section*{References}
\begin{enumerate}
  \item MedlinePlus. Amputation. U.S. National Library of Medicine; 2025. Available from: \url{https://medlineplus.gov/amputation.html}
  \item Townsend CM, Beauchamp RD, Evers BM, Mattox KL, editors. Sabiston Textbook of Surgery: The Biological Basis of Modern Surgical Practice. 22nd ed. Elsevier; 2025.
  \item American Academy of Orthopaedic Surgeons (AAOS). Amputation. Available from: \url{https://www.aaos.org/conditions/amputation/}
  \item Society for Vascular Surgery (SVS). Clinical Practice Guidelines for the Management of Patients with Critical Limb Ischemia. J Vasc Surg. 2025; (forthcoming).
\end{enumerate}

\clearpage
